\section{Conclusion}

This paper developed a Euclidean support-coordinate description of scale
and shape for a fixed-normal family based on the regular dodecahedron.

The six pairs of opposite face-normal directions define the support
register
\[
\mathbf h=(h_1,\ldots,h_6)\in\mathbb R_{>0}^6.
\]

Its realization is
\[
P(\mathbf h)
=
\left\{
\mathbf x\in\mathbb E^3:
|\mathbf n_i\cdot\mathbf x|\le h_i
\text{ for }i=1,\ldots,6
\right\}.
\]

Within the active fixed-normal family, these six coordinates uniquely
determine the polytope.

The support space admits the canonical orthogonal decomposition
\[
\mathbb R^6=U\oplus H,
\]
where
\[
U=\operatorname{span}\{\mathbf 1\}
\]
and
\[
H=
\left\{
\mathbf r\in\mathbb R^6:
\sum_{i=1}^{6}r_i=0
\right\}.
\]

Accordingly,
\[
\mathbf h
=
\bar h\,\mathbf 1+\mathbf r,
\qquad
\bar h
=
\frac{1}{6}\sum_{i=1}^{6}h_i.
\]

The geometric meaning of this split is exact.

For every \(c>0\),
\[
P(c\mathbf h)=cP(\mathbf h).
\]

Thus the positive uniform ray gives the regular similarity family. While all
defining planes remain active facets, a realization homothetic to the
regular reference dodecahedron must have equal support values. The uniform
ray is therefore the unique regular scale direction.

Within the dodecahedral support chamber,
\[
\mathbf r=\mathbf 0
\]
if and only if the realization is regular. Every nonzero residual produces
a non-homothetic support variation.

The decomposition is preserved by the rotational symmetry group. Its
orthogonal projections are
\[
P_U=\frac{1}{6}J
\]
and
\[
P_H=I-\frac{1}{6}J.
\]

Every nonregular support state has the unique form
\[
\mathbf h
=
\bar h\,\mathbf 1
+
\sigma\widehat{\mathbf r},
\]
where
\[
\sigma=\|\mathbf r\|>0
\]
and
\[
\widehat{\mathbf r}\in S^4\subset H.
\]

Equivalently,
\[
\mathbf h
=
\bar h
\left(
\mathbf 1+\eta\widehat{\mathbf r}
\right),
\qquad
\eta=\frac{\sigma}{\bar h}.
\]

The admitted family therefore carries one scale coordinate and five shape
degrees of freedom, organized by magnitude and direction.

The result is bounded to origin-centered, centrally symmetric realizations
with the regular dodecahedral normal directions fixed and the dodecahedral
incidence structure preserved. It does not classify arbitrary Euclidean
dodecahedra or arbitrary deformations of the regular solid.

No mechanics, elasticity, constitutive law, force, energy, dynamics,
continuum model, physical expansion, gravitation, or cosmological
interpretation has been introduced.

The main conclusion is classical:

\begin{quote}
For the fixed-normal dodecahedral support family, Euclidean convex geometry
and finite-dimensional real linear algebra are sufficient to separate
regular scale from relative shape.
\end{quote}

The construction does not require mathematics to be extended. It requires
the available mathematics to be used carefully and carried to its natural
boundary.
